\section{Mechanism}
\label{sec:method}

\subsection{Setting}

Let $\theta_t$ denote the shared model at distributed round $t$. There are $N$
workers and each worker performs $K$ local optimization steps before the server
updates the shared state. Throughout this paper, $K$ is the number of local
steps between aggregation events. A standard local-SGD-style update collects
worker deltas $u_{t,i}$ and applies a mean or optimizer transform to obtain
$\theta_{t+1}$.

\method changes the round by introducing an arm set $\arms$. An arm can be a
worker-side hyperparameter, such as a local learning-rate multiplier, or a
server-side control, such as an outer learning rate and momentum pair applied
to a common pool of worker updates. Each arm $h \in \arms$ induces a candidate
next state
\begin{equation}
\theta_{t+1}^{(h)} = \theta_t + A_t(U_t, h),
\end{equation}
where $U_t$ is the set of worker updates available at the server and $A_t$ is
the arm-specific update map. A probe batch $\probe$ estimates a loss
$\ell_t(h)=\loss(\theta_{t+1}^{(h)}; \probe)$. The server converts these losses
to weights or a one-hot selection and applies a single shared update.

\subsection{Worker-arm aggregation}

In the worker-arm version, workers or worker groups receive different local
optimizer settings. For example, groups may share the same starting state
$\theta_t$ but run $K$ local steps with learning-rate multipliers
$h \in \{0.25,0.5,1,2,4\}$. The returned deltas are grouped by arm, scored by
probing the resulting candidate next states, and then combined:
\begin{equation}
\theta_{t+1} = \theta_t + \sum_{h \in \arms} w_t(h)\,\bar{u}_t(h),
\end{equation}
where $\bar{u}_t(h)$ is the candidate update for arm $h$ and $w_t$ is derived
from probe scores.

Worker-arm aggregation is the most direct version of the idea, but it has a
cost tradeoff: with a fixed number of workers, each arm receives fewer workers.
This is why our experiments always compare to fixed-compute external HPO and
to shuffled-probe controls. A positive result must come from useful arm
information, not from giving the method more total update work.

\subsection{Server-control aggregation}

In the server-control version, workers can all produce a common update pool.
The server then constructs several candidate outer updates from the same pool.
For example, a candidate may use outer step size $\alpha$ and momentum
$\beta$:
\begin{equation}
v_{t+1}^{(h)} = \beta_h v_t + \bar{u}_t,\qquad
\theta_{t+1}^{(h)} = \theta_t + \alpha_h v_{t+1}^{(h)}.
\end{equation}
The probe scores $\theta_{t+1}^{(h)}$ for each $h=(\alpha_h,\beta_h)$ and the
server applies the selected or weighted candidate. This version directly
addresses the concern that internal HPO simply starves each parameter setting
of workers: the arms can be generated after the worker updates have already
been collected.

\subsection{Probe-to-weight rules}

The simplest rule is one-hot selection,
\begin{equation}
h_t^\star \in \arg\min_{h \in \arms} \ell_t(h),\qquad
w_t(h)=\mathbf{1}\{h=h_t^\star\}.
\end{equation}
We also use conservative soft weights in some local-SGD ledgers:
\begin{equation}
q_t(h) \propto \exp\left(-\frac{\ell_t(h)-\bar{\ell}_t}{\tau_t}\right),
\qquad
w_t(h) = (1-\lambda)\pi_t(h) + \lambda q_t(h),
\end{equation}
where $\pi_t$ is a prior or default mixture, $\tau_t$ scales with the observed
probe dispersion, and $\lambda$ shrinks the probe response. One-hot selection
is the limiting case used in the adaptive HPO experiments when the probe is
trusted to choose among a small fixed arm set.

\subsection{Cost accounting}

The probe must be cheap enough for internal HPO to be meaningful. If $G$
candidate arms are evaluated on $S$ probe slices of size $B_p$, and the
training round performs $N K$ local worker steps with training batch size $B$,
the forward-pass probe overhead is approximately
\begin{equation}
\mathrm{overhead}
  \approx
  \frac{G S B_p}{3 N K B},
\end{equation}
where the factor $3$ is the standard rough forward/backward cost ratio used in
our ledgers. In server-control arms, the worker-update cost is shared across
candidates. In the adaptive shift audits, the reported overhead is the measured
probe-forward ledger, with a low-probe rerun at 8.33\%.
